\begin{figure}[p]
    
    \begin{minipage}{\textwidth}
        \begin{mdframed}
            \label{alg:fulloneserv}
            \begin{center}
                {\bf Single-Server Scheme for $Q=\sqrt{n}/2$ Queries
                    \footnote{We first present the scheme supporting distinct and random queries. As mentioned, these restrictions can be removed by applying PRP and local caching.}
                }
            \end{center}
            \parhead{Notation.}
            $\kappa$ denotes a {\it statistical} security parameter, $\lambda$
            denotes a computational security parameter. We use $\alpha(\kappa)$
            to denote an arbitrarily small super-constant function.
            
            
            \parhead{Preprocessing.} 
            %\ \\[-6pt]
            \begin{itemize}[leftmargin=3mm,itemsep=1pt]
                \item 
                \Client samples $M_1=2\sqrt{n}\log \kappa \cdot \alpha(\kappa)$
                master \ppps keys denoted $\psk_1, \dots, \psk_{M_1} \in \{0,
                1\}^\lambda$.
                %, where $M_1=\sqrt{n}\log \kappa \cdot \alpha(\kappa)$. , such
                %that $\alpha(n)=\omega(1)$. \elaine{TODO: fix params}
                Initialize the parities $p_1,\dots,p_{M_1}$ to zeros.
                
                %\elaine{TODO: notation inconsistent with earlier text}
                
                \item 
                \Client downloads the whole DB from the server in a streaming
                way: when the client has the $j$-th chunk
                $\DB[j\sqrt{n}:(j+1)\sqrt{n}]$:
                \begin{itemize}
                    \item {\it Update primary table}: for $i \in [M_1]$, let
                    $p_i \get p_i \oplus \DB[\Set(\psk_i)[j]]$.
                    \item {\it Store replacement entries}: sample and store
                    $M_2=3\log \kappa \cdot \alpha(\kappa)$ % \log \kappa\cdot
                    % \alpha(\kappa)$ 
                    tuples of the form $(r, \DB[r])$ where $r$ is a random index
                    from the $j$-th chunk.
                    %independent random indices from $[i\sqrt{n}:(i+1)\sqrt{n}]$
                    %with the associated database values.
                    \item Delete $\DB[j \sqrt{n}:(j +1)\sqrt{n}]$ from the local
                    storage.
                \end{itemize}
                
                \item At this moment, let $T := \{(\psk_i, p_i )\}_{i \in
                [M_1]}$ denote the client's {\it hint table}. Mark all the hints
                as ``good''.
            \end{itemize}
            
            
            
            
            \parhead{Query for index $x \in \{0, 1, \ldots, n-1\}$.}
            %When a query $x \in \{0, 1, \ldots, n-1\}$ arrives, \Client
            %performs the following: 
%            \begin{enumerate}[leftmargin=5mm,itemsep=1pt]
%                \item 
%                {\bf Query:}
                \begin{enumerate}[leftmargin=5mm,itemsep=1pt]
\item 
${\sf Client}$:
 For each matched entry $(\sk_i, p_i)$ such that $x\in \Set(\sk_i)$
in the hint table, do the following unless there are already $M_3 = 3 \log \kappa \cdot \alpha$ matched entries: 
\begin{itemize}[leftmargin=4mm]
\item 
For the first good (i.e., non-broken) matched entry,  
find the first unconsumed replacement entry $(r, \DB[r])$ 
for $\chunk(x)$. \footnote{If not found, treat it as the otherwise case.}
\item 
Otherwise, sample a random index %$r \in \{0, \ldots, \sqrt{n}-1\}$.
$r$ in $\chunk(x)$.
\item 
$(\sk', i^*)\leftarrow {\sf Program}(\sk_i, {\sf chunk}(x), r \text{ mod } \sqrt{n})$.
\item 
Send $\sk'$ to the server, and receive  
$\{\beta_i\}_{i \in \{0, \ldots, n^{1/4}-1\}}$ from the server.
\item 
For the first good matched entry, save the answer 
$p_i \oplus \beta_{i^*} \oplus \DB[r]$. %where $\DB[r]$ is retrieved
%from the replacement entry consumed.
\item 
Sample a fresh 
\ppps key $\sk_{\rm new}$ subject to $x \in \Set(\sk)$, and 
replace the 
consumed entry $(\sk_i, p_i)$ with 
$(\sk_{\rm new}, 0)$
and mark the entry as {\it broken}.
\end{itemize}

\item 
${\sf Client}$:
If less than $M_3$ keys are sent in the previous step, send more dummy programmed keys to the server until there are $M_3$ keys sent~\footnote{The dummy key is constructed as sampling a random \ppps key $\sk$ subject to $x\in \Set(\sk)$ and call $\sk', \_ \leftarrow {\sf Program}(\sk, \chunk(x), \delta')$.}.
%repeat the following $M_3 - {\sf cnt}$ times:
%sample a random \ppps key $\sk$ subject to $x \in \Set(\sk)$, 
%call 
%$\sk', \_ \leftarrow {\sf Program}(\sk, \chunk(x), \delta')$
%where $\delta'$ is 
%sampled randomly from $\{0, 1,\ldots, \sqrt{n}-1\}$,
%and send $\sk'$ to the server. Ignore the answers received.

\item 
${\sf Client}$: Output the saved answer.  
If no answer was saved, output $0$.
\item 
${\sf Server}$: 
for each $\sk'$ received, 
let $S_0, \ldots, S_{n^{1/4}-1} \leftarrow {\sf ListDecode}(\sk')$. 
For each $i \in \{0, \ldots, n^{1/4}-1\}$, 
send the xor-sum $\oplus_{j \in S_i}\DB[j]$ to the  
client\footnote{We use the normal representation of the set $S_i$ and not
the offset representation.}. 
\ignore{
\item 
${\sf Client}$: 
Let $g$ be the index of the first 
matched entry that is not marked broken --- if no such $g$ is found,
simply output $0$. 
Let
$\{\beta_i\}_{i \in [n^{1/4}]}$ 
be the answers received from the server
corresponding to the $g$-th programmed key $\sk'$ sent to the server.  
Output 
$\beta_i \oplus \DB[r]$
}


\ignore{
                    \Client finds all the matched hints
                    $(\sk_1,p_1),\dots,(\sk_m,p_m)$ in the hint table  such that
                    $x\in \Set(\sk_j)$ for all $j\in[m]$.
                    
                    \item 
                    If there are less than $M_3=6(\ln \kappa \alpha(\kappa))^2$
                    hints, pad them with random broken hints. If there are more
                    than $M_3$ hints, sample $M_3$ hints from them. Randomly
                    shuffle the order of these hints\footnote{We add this random
                    shuffle to ensure that the order of the hints do not reveal
                    any information}.
                    %\item If there are more than $M_3$ matched hints, randomly
                    %choose $M_3$ hints among them. (see \Cref{eqn:setfromkey}
                    %for the definition of $\Set(\psk_j, x')$ where $x'$ is
                    %possibly $\bot$). 
                    
                    
                    % \item For the one of the randomly chosen good matched hint
                    % $(\sk_j, p_j)$\footnote{If there are no good matched hint,
                    % just mark the answer as 0.}:
                    % \begin{itemize}
                    %     \item \Client finds the first unconsumed replacement
                    %     entries $(r, \DB[r])$ from $x$'s chunk;
                    %     \item Execute $\SingleQuery(x, \sk_j, p_j, r, \DB[r])$
                    %     and mark the answer for the query.
                    % \end{itemize} 
                    
                    % \item Simultaneously, for other matched hints $(\sk_j,
                    % p_j)$: 
                    % \begin{itemize}
                    %     \item \Client executes $\SingleQuery(x, \sk_j, 0, r',
                    %     0)$ with a random index $r'$ from $x$'s chunk and
                    %     ignores the answer.
                    % \end{itemize}
                    \item For each hint, do the following:
            \begin{itemize}
                \item if the hint is marked good, find the first unconsumed replacement entry $(r, \DB[r])$ from $x$'s chunk;
                and execute $\SingleQuery(x, \sk_j, p_j,r,\DB[r])$ and store the answer for the query.
                \item Otherwise, execute $\SingleQuery(x, \sk_j, 0,0,r')$ with a random index $r'$ from $x$'s chunk  and ignore the answer.
            \end{itemize}
                
            \item Return the answer of any of the good hints. If none of the hints were good, just return $0$.
}
                    
%                \end{enumerate}
                
\ignore{
                \item 
                {\bf Refresh:}
                \begin{itemize}
                    \item \Client replaces all the keys for the matched hints in
                    the hint table with $\sk'_1,\dots,\sk'_m$, such that for all
                    $j$, $\sk'_j$ is rejection-sampled conditioned on $x\in
                    \Set(\sk'_j)$. Also mark the parities as 0 and mark those
                    hints as ``broken''.
                \end{itemize}
       }         
            \end{enumerate}
        \end{mdframed}
    \end{minipage}
    \caption{
Our single-server pre-processing PIR scheme.
%Sublinear single server preprocessing PIR protocol introduced in
%    Section ~\labelcref{sec:single-server}. The protocol uses $\SingleQuery$ defined
%    in Fig.~\labelcref{fig:single-query} as a sub-routine.
}
    \label{fig:1server-new}
\end{figure}


\section{Our Single-Server PIR Scheme}\label{sec:one-server}


\subsection{Construction}

\paragraph{Notation.} For $x\in \{0,1,\ldots,n-1\}$, we define $\chunk(x):=\lfloor x/n^{1/2} \rfloor$ and $\superblock(x):=\lfloor \chunk(x)/n^{1/4} \rfloor$. We assume $(\Gen,\Set,\Program,\Decode)$ is a \ppps scheme over the distribution $\mcal{D}_n$ as described in Section~\labelcref{sec:ppps}.

\paragraph{Intuition.} The major differences between our single-server scheme and the two-server scheme are summarized below.

\begin{itemize}[leftmargin=*]
    \item \textit{Preprocessing.} The two server scheme allows the client to do preprocessing with one server and do online queries with another server. 
    Our single-server scheme uses the technique from Piano~\cite{piano} such that the client makes a streaming pass over the whole database (retrieving from the only server) and runs the preprocessing locally. 
    
    \item \textit{Query and Refresh.}
    The two-server scheme allows the client to replace a consumed set with a new set on-the-fly, because the client can query another server for the new parity.
    Instead, in the single-serve scheme, we use a new \underline{broken hint} idea. 
    The client still generates a new set after the query, but it only marks the new set as ``broken hint'' since the parity is unknown. 
    To ensure correctness, the client now  \textit{uses all the matched sets} given one query, and as long as there is one non-broken hint, the answer can be computed correctly.
\end{itemize}


\paragraph{Detailed algorithm for bounded, random queries.}
%We describe the detailed construction
In \Cref{fig:1server-new}, we describe
our algorithm which supports  
$Q = \sqrt{n}/2$ random and distinct queries.
It is well-known how to upgrade such an algorithm
to support an {\it unbounded} number of 
{\it arbitrary} queries~\cite{pir-ec23}.
For completeness, we briefly describe the upgrade shortly after.

\paragraph{Efficiency.}
Observe that although the list decoding
produces $n^{1/4}$ candidate sets
each of size $\sqrt{n}$, all $n^{1/4}$ 
sets can actually be represented using 
only $O(\sqrt{n})$ 
space. Further, computing the parities 
of all sets 
takes only $O(\sqrt{n})$ time.
This is because all $n^{1/4}$ sets are derived
from some common set $S$, but replacing 
offsets within each of the $n^{1/4}$ superblocks
with another random vector $\delta_0, \ldots, \delta_{n^{1/4}-1}$.
We give a full efficiency analysis
in the proof of Theorem~\labelcref{thm:eff1srv}.


\paragraph{Supporting unbounded, arbitrary queries.}
% \elaine{TODO: add}
We can easily get rid of the distinct query
assumption in the following way: we can require the client to store a local cache of size $Q$ for caching the
most recent $Q$ queries. If the client wants a repeated query, it can lookup
in the cache and make a dummy query. 

Further, we can assume that the queries are random without loss of generality as follows: we can let the client and the server agree on a pseudorandom permutation (PRP)~\cite{prp1,prp2} upfront and the server can permute the database according to the PRP. Another option is let the server
build the database as a key-value storage and use a cuckoo hash
table~\cite{cuckoo1,cuckoo2} directly based on a PRF to locate the queries.
Notice that, in both implementations, the client can still make queries
adaptively depending on the real query sequence and the responses, which is
sufficient for practical usage. Then, as long as the client makes the
queries independent of the randomness of the PRP/PRF, those queries can be
considered as uniformly random. This assumption is only needed for the
correctness.

Lastly, we can remove the bounded $Q$ query assumption as follows: the straightforward way is that once the client finishes a window of $Q$ queries, the client and the server reruns the preprocessing phase again, using fresh randomness. 
The drawback is that the client has to wait a long time before starting the next window.
As previous work pointed out~\cite{piano,pir-ec23}, we can easily avoid
this drawback 
through a simple pipelining trick, by spreading
the preprocessing work of the next $Q$ window
over the current $Q$ window of queries. 
\ignore{
We can do something relatively better using a simple pipelining trick, similar to that suggested in Piano~\cite{piano}. 
The scheme first runs a preprocessing phase for the first $Q$ queries. 
Then, during the query phase, the client interacts with the preprocessing server and gradually finishes the preprocessing for the next $Q$ queries. 
Essentially, the client and the server can split the work (including computation and communication) of the preprocessing into $Q$ pieces.
Whenever the client finishes a query, the client and the server will finish the next piece of work about preprocessing.
Once the client finishes the $Q$ queries, it can immediately start the next window of $Q$ queries because the preprocessing for the next window is done. 
The main advantage is that the client never waits long time after the first preprocessing phase.
The downside is: 1) in the two server setting, it requires the preprocessing server to be always online; 2) in the single-server setting, the online communication will be dominated by the preprocessing parts and become $O(\sqrt{n})$, instead of $\widetilde{O}_\lambda(n^{1/4})$.
}

\ignore{
    The client only stores $M_1=2\sqrt{n}\log \kappa\alpha(\kappa)$
    $\lambda$-bit keys and store in total $\sqrt{n}\cdot M_2=3\sqrt{n}\log
    \kappa\alpha(\kappa)$ index-value pairs. So the storage is
    $O(\lambda\sqrt{n}\log \kappa \alpha(\kappa))$.
    
    The pre-processing phase's performance bounds follow straightforward by the algorithm descriptions.
    
    For the query phase, the client first enumerates all $M_1$ hints to find
    $x$, which takes $O(\lambda\sqrt{n}\log \kappa\alpha(\kappa))$ time. For all
    the $M_3=\Theta((\log \kappa \alpha(\kappa))^2)$ found hints, the
    $\SingleQuery$ algorithm takes $O(\lambda n^{1/4})$ client computation and
    $O(\lambda \sqrt{n})$ server computation, while the communication cost is $O(\lambda
    n^{1/4})$. During the refresh phase, we consider the expected computation time: each hint in the table will have $1/\sqrt{n}$ probability to be resampled and each resample takes $O(\lambda\sqrt{n})$ expected time to finish. 
    So the total expected time is $O(\lambda\sqrt{n}\log \kappa
    \alpha(\kappa))$  per query. The server time is
    $O(\lambda\sqrt{n}(\log \kappa \alpha(\kappa))^2)$ per query. The online
    communication per query is $O(\lambda n^{1/4}(\log \kappa
    \alpha(\kappa))^2)$.
}




\begin{theorem}
    \label{thm:main-1server} %Let $\lambda,\kappa\in \mathbb{N}$.
    \label{thm:eff1srv}
    Let $\alpha(\kappa)$ be any super-constant function. 
Suppose that $\PRF_1,\PRF_2$ are secure pseudorandom functions, and 
$n$ is bounded by $\poly(\lambda)$ and
    $\poly(\kappa)$. 
The single-server scheme in~\Cref{fig:1server-new} that supports
    $\sqrt{n}/2$ random, distinct queries is private, and correct with probability
    $1-\negl(\lambda)-\negl(\kappa)$ for some negligible function $\negl(\cdot)$.
Further, it achieves the following performance
    bounds:
    \begin{itemize}
        \item $O_\lambda(\sqrt{n}\log \kappa \cdot \alpha)$ client storage and
        no additional server storage;
        \item \textbf{Pre-processing Phase}:
        \begin{itemize}
\item $O(n)$ server time and $O_\lambda(n\log \kappa \cdot \alpha)$ 
client time;
            \item {$O(n)$ communication};
        \end{itemize}
        \item \textbf{Query Phase}:
        \begin{itemize}
            \item $O_\lambda(\sqrt{n}\log \kappa \cdot \alpha)$ expected client
            time and $O_\lambda(\sqrt{n}\log \kappa \cdot \alpha)$ server
            time per query; 
            \item $O_\lambda(n^{1/4}\log \kappa\cdot \alpha)$
            communication per query.
        \end{itemize}
    \end{itemize} 
Therefore, the amortized 
online communication per query is $O_\lambda(n^{1/4} \log \kappa 
\cdot \alpha)$, 
the amortized offline communication per query 
is $O(\sqrt{n})$, the amortized
client and server computation per query is
$O_\lambda(\sqrt{n}\log \kappa \cdot \alpha)$.
\end{theorem}

\begin{proof}
    We defer the privacy and the correctness proof to
\Cref{sec:1server-priv}
and 
\Cref{sec:1server-correct}
respectively.
%~\Cref{sec:proof}.
    The client only stores $M_1=2\sqrt{n}\log \kappa\alpha$
    $\lambda$-bit keys and store in total $\sqrt{n}\cdot M_2=3\sqrt{n}\log
    \kappa\cdot \alpha$ index-value pairs. So the storage is
    $O_\lambda(\sqrt{n}\log \kappa \cdot \alpha)$.
    The pre-processing phase's performance bounds follow straightforward by the algorithm descriptions.
    
    For the query phase, the client first enumerates all $M_1$ hints to find
    $x$, which takes $O_\lambda(\sqrt{n}\log \kappa\cdot \alpha)$ time. For all
    the $M_3=\Theta(\log \kappa \cdot \alpha)$ found hints, the
    $\Program$ algorithm takes $O_{\lambda}(n^{1/4})$ client computation. For the server, during the query phase, the client sends $M_3$ programmed \ppps key to the server, the size of which is $O_\lambda(n^{1/4})$. The server sends back the xor-sum of $n^{1/4}$ candidate sets for each key after running $\ListDecode$. Even though each candidate set has size $n^{1/2}$, all the candidate sets have a succinct representation of size $n^{1/2}$ and server can compute this representation in time $O_\lambda(n^{1/2})$. Further, as observed in Theorem~\labelcref{thm:eff2srv}, due to the structure of the candidate sets, the sever can compute all the $n^{1/4}$ xor-sums in time only $O(n^{1/2})$. Hence, the server's running time %for each hint is
    $O_{\lambda}(\sqrt{n})$. 
    
    The client needs to find $M_3$ matched hints, compute $O(\log \kappa \cdot \alpha)$ xor operations during each query and for the matched hints, it needs to sample fresh \ppps key $\sk_\new$ subject to $x\in \Set(\sk_\new)$. The client's computation time is dominated by this sampling step. We consider the expected computation time: each key in the hint table will have $1/\sqrt{n}$ probability to be replaced in this and each sampling takes $O_\lambda(\sqrt{n})$ expected time to finish using Lemma~\labelcref{lem:hintone} and pseudorandomness of \ppps, the expected number of keys checked until a key $\sk$ such that $\Set(\sk)$ contains the current query is found is $O(\sqrt{n})$.
    So the total expected time for the query phase is $O_\lambda(\sqrt{n}\log \kappa
    \alpha)$  per query. The server time is
    $O_\lambda(\sqrt{n}(\log \kappa \cdot\alpha))$ per query. The online
    communication per query is $O_\lambda( n^{1/4}(\log \kappa
    \cdot \alpha))$.
\end{proof}


\subsection{Privacy Proof}
\label{sec:1server-priv}

Suppose that the underlying \ppps scheme satisfies
private programmability.
Below, we prove the privacy of our single-server PIR scheme.

In the pre-processing 
phase, the server observes a single scan over the database,
and thus no information is leaked.
The rest of the proof will therefore
focus on the query phase.

\begin{lemma}\label{lem:hintone}
Recall that in each time step $t$, the adversary \algA 
adaptively chooses a query $x_t \in \{0, 1, \ldots, n-1\}$ for the client.
%The following statements are true:
%\begin{enumerate}[leftmargin=5mm]
%\item 
At the end of each time step $t$, 
the client's hint table  
is distributed 
as a table of size $M_1$
where each entry is a freshly sampled \ppps key, even when
conditioned on $\algA$'s view so far.  
\ignore{
\item 
In each time step $t$, 
the programmed key $\sk'$ 
sent to $\algA$ 
is the same as sampling a random \ppps key 
}
%\end{enumerate}
\label{lem:hintfresh}
\end{lemma}
\begin{proof}
Suppose the above statement holds at the end of time step $t-1$.
We prove that it still holds 
at the end of time step $t$.
Since the hint table is distributed as a fresh randomly sampled
table even when conditioned on $\algA$'s view
at the end of $t-1$, we may henceforth assume
an arbitrary fixed query $x_t$.
The distribution of the hint table 
before the $t$-th query can be equivalently rewritten as:
\begin{itemize}[leftmargin=5mm]
\item 
First, sample the indices of the entries (henceforth denoted $I$)
that contain the query $x_t$. Specifically, each  
$i \in [M_1]$ is chosen into the set $I$ 
independently with probability $1/\sqrt{n}$.
\item 
For each $i \notin I$, sample a random \ppps
key subject to not containing $x_t$.
\item 
For each $i \in I$, sample a random \ppps
key subject to containing $x_t$.
\end{itemize}

Using the above interpretation, it is easy to see 
that the distribution of the hint table 
after the $t$-th query is unaltered.
\end{proof}


\paragraph{Simulator construction and the ${\sf Ideal}$ experiment.}
Consider the following simulator construction which does not make
use of the queries: 
in every time step $t$, call the simulator $\Sim$ of the \ppps
scheme, and let the output be $\sk'$. 
Send $\sk'$ to the server. 

\paragraph{Indistinguishability of ${\sf Real}$ and 
${\sf Ideal}$.}
We now prove the indistinguishability of the ${\sf Real}$ and
${\sf Ideal}$
assuming the 
private programmability of the underlying \ppps scheme.

First, due to \Cref{lem:hintfresh}, 
we can equivalently rewrite the ${\sf Real}$ experiment as follows:
at the end of each time step, resample the entire hint table
freshly at random before continuing to answer more queries.
As a result, the %distribution of the $\sk'$ received by 
messages sent to 
\algA in each time step $t$
is distributed as 
\begin{itemize}[leftmargin=5mm]
\item[]
Repeat $M_3$ times: 
\begin{itemize}[leftmargin=5mm]
\item 
sample a \ppps key $\sk$
 at random subject to containing the query $x_t$;
sample $\delta$ at random
from $\{0, \ldots, \sqrt{n}-1\}$.
\item 
call $\sk', \_ \leftarrow {\sf Program}(\sk, \chunk(x_t), \delta)$; 
\item 
send $\sk'$ to \algA.
\end{itemize}
\end{itemize}
One way to see this is to think of the distribution of the table
as the equivalent distribution
in the proof of \Cref{lem:hintfresh}.

Therefore, the rest of the proof follows
due to a straightforward hybrid argument
where we replace
the programmed keys (denoted $\sk'$ earlier) 
sent to the server in all time steps one by one
with a simulated key, 
relying on the private programmability of the underlying \ppps.

\subsection{Correctness Proof}
\label{sec:1server-correct}


    %Recall that in our full scheme, the server will first sample a pseudorandom permutation~(PRP) to permute the database upfront and the client will download the key from the server. 
    %Replacing the PRP with a true random permutation only affects the failure probability by a negligible amount, $\negl(\lambda)$.
    %We also assume that the client does not make any duplicative queries for those $Q$ queries.
    %Therefore, taking the randomness of the permutation, we can view all $Q$ queries are randomly sampled from $\{0,1,\dots,n-1\}$ without replacement.
    
%    We assume the client uses a true random oracle to sample the sets, instead of a PRF. 
%    Due to the pseudorandomness of the PRF, this assumption will not affect the failure probability by a negligible amount, $\negl(\lambda)$. 
For the correctness analysis, 
we may assume that every set
is sampled independently from $\mcal{D}_n$. 
Due to the pseudorandomness property of the \ppps,
this will only affect the correctness probability 
by a negligible amount.
   
Recall that we have a window of $Q = \sqrt{n}/2$
random, distinct queries.
    There are only two 
bad events that can cause correct failure: 
1) the client cannot find a good hint that contains the query index; 2) the client runs out of replacements in a chunk.

    We first analyze the second bad event, i.e., depleting replacement entries
For every query, at most one replacement
entry is consumed. 
Therefore, the second bad event  
only happens when the client makes more than $M_2$ queries in one chunk. 
The analysis is the same 
as analysis of depleting replacement entries
in the proof of correctness 
for our 2-server scheme (see \Cref{sec:2server-correct}).
\ignore{
Since the client is making $\sqrt{n}\ln \kappa \alpha(\kappa)$ queries and there are $\sqrt{n}$ groups, we can use a standard balls-into-bins argument.  
    For $t\in [Q], i\in \{0,1,\dots,\sqrt{n}-1\}$, define the random variables $Y_{t,i}\in \{0,1\}$ such that $Y_{t,i}=1$ if and only if the $t$-th query locates in the $i$-th chunk.
    Denote $X_i=Y_{1,i}+\dots+Y_{t,i}$ be the number of queries located in the $i$-th chunk. 
    We know $\E[Y_{t,i}]=1/\sqrt{n}$ and $\E[X_i]=\ln \kappa \alpha(\kappa)$.
    Notice that we are taking the randomness of the permutation and the queries do not have duplication, so $Y_{1,1},\dots,Y_{Q,1}$ are negatively correlated.     
    With the Chernoff bound for negatively correlated variables, we know that 
    \begin{align*}
        & \Pr[X_1 \ge (1+2)\ln \kappa \alpha(\kappa) ] \\ 
        \le & \exp \left(\frac{-2^2}{2+2}\ln
        \kappa\alpha(\kappa)\right) = \kappa^{-\Theta(\alpha(\kappa))}.
    \end{align*}
    
    Taking the union bound over all $\sqrt{n}$ chunks, the failure probability is bounded by $\sqrt{n}\cdot \kappa^{-\Theta(\alpha(\kappa))}$, which is a negligible function of $\kappa$.% since $n=\poly(\kappa)$.
}
    
    The first bad event, i.e., no matched good hint, 
can only arise from the following events: 
1) there are no good hints left that match the query; 
2) there are more than $M_3$ matched hints 
but the first $M_3$ matched hints are all broken. 
    
    Fix a sequence of query $x_1,\dots,x_m$ and consider the 
error probability for $x_m$. 
    Consider the initial hint table
in which each entry represents a random set sampled from $\mcal{D}_n$. 
%, whose generation process is independent of $x_1,\dots,x_m$. %due to the assumption of random queries.
    If a hint in the initial hint table contains $x_m$ and does not contain $x_1,\dots,x_{m-1}$, this hint will remain good until the query for $x_m$. We have that, for any hint
    \begin{displaymath}
        \Prob{\text{ the hint contains }x_m}=1/\sqrt{n}\;.
    \end{displaymath}
    % We know that each independent hint contains $x$ with probability $1/\sqrt{n}$. 
Further, conditioned on a hint containing $x_m$, if $x_i$ and $x_m$ are not in the same chunk, then the hint contains $x_i$ with probability $1/\sqrt{n}$ because each chunk has its own randomness.
Moreover, if $x_i$ and $x_m$ are in the same chunk, this hint definitely will not contain $x_i$.
    Hence, we have that 
    \begin{align*}
        &\CondProb{\text{ hint is broken}}{\text{ the hint contains }x_m}\\
        \leq& \CondProb{\exists i\in \{1,\ldots,m-1\}\text{ hint contains }x_i}{\text{ the hint contains }x_m}\\
%        \leq &\Prob{\exists i\in \{1,\ldots,m-1\}\text{ hint contains }x_i }
\leq & (m-1)/\sqrt{n}\leq Q/\sqrt{n}\leq 1/2\;.
    \end{align*}
%    The second inequality above follows because a hint cannot be broken if it did not contain any earlier $x_i$. The third inequality above follows because the events that $x_i$'s are contained in a hint are independent. 
%    The last inequality follows because $Q\le \sqrt{n}/2$.
    Therefore, 
    \begin{align*}
        &\Prob{\text{ hint is good and contains }x_m}\\
        \geq &\CondProb{\text{ hint is good}}{\text{ the hint contains }x_m}\cdot \Prob{\text{ the hint contains }x_m} \geq 1/(2\sqrt{n})\;.
    \end{align*}
    Then, the probability that 
there is no good hint matching $x_m$ in the table is 
at most 
    \begin{align*}
        & \left(1-\frac{1}{2\sqrt{n}}\right)^{M_1}  = \left(1-\frac{1}{2\sqrt{n}}\right)^{2\sqrt{n}\ln \kappa\cdot \alpha(\kappa)}  \\
        \le & \left(1/e\right)^{\ln \kappa \alpha(\kappa)}  = \kappa^{-\alpha(\kappa)}.
    \end{align*}
Now let us argue that for query $x_m$, the probability of more than $M_3$ hints being matched is small. 
    Due to the Lemma~\ref{lem:hintfresh}, 
we may assume that at the beginning of each query,
the hint table contains freshly and independently chosen sets.
%the local hints will be identically distributed as freshly-random and each hint will contain the query with probability $1/\sqrt{n}$. 
Each set sampled from $\mcal{D}_n$ 
contains the query with probability $1/\sqrt{n}$.
The expected number of hints that match
the query is therefore $M_1/\sqrt{n} = 2 \log \kappa  \cdot \alpha$.
    Using the Chernoff bound, we have that the probability of more than more than $M_3 = 3 \log \kappa \cdot \alpha$ hints being matched is bounded by
a negligibly small function in $\kappa$.
\ignore{
    \begin{align*}
     & \Pr[\#\textsf{matched hints} \ge (1+2) (M_1/\sqrt{n})\ln \kappa \alpha(\kappa) ] \\ 
     \le & \exp \left(\frac{-2^2}{2+2}\ln
     \kappa\alpha(\kappa)\right) = \kappa^{-\Theta(\alpha(\kappa))}.
    \end{align*}   
}

Finally, we can apply a union bound
over all $\sqrt{n}/2$ queries and conclude that 
the probability of 
the first bad event (i.e., no matched good hint)
ever happening is negligibly small in $\kappa$.
\ignore{
With the union bound, for all $\sqrt{n}/2$ queries, the probability of any query cannot find a set is bounded by $\frac{1}{2}\sqrt{n}\cdot 
    (\kappa^{-\alpha(\kappa)} + \kappa^{-\Theta(\alpha(\kappa))})$, which is negligible in $\kappa$ since $n$ is bounded by $\poly(\kappa)$.
    
    
    Then, there is some negligible function $\negl(\cdot)$ that all the queries are answered correctly with probability at least $1-\negl(\lambda)-\negl(\kappa)$.
    %This proves the statement.
   } 

